\chapter{Tổng quan}
\label{chap:overview}

\section{Giới thiệu đề tài}

Phần này dùng để mô tả bối cảnh của đề tài, vì sao bài toán này quan trọng,
và nhóm đang giải quyết nhu cầu nào trong thực tế. Nên viết ngắn gọn, rõ ràng,
và ưu tiên góc nhìn của người đọc chưa biết gì về đồ án.

\section{Mục tiêu}

Có thể chia mục tiêu thành hai nhóm:

\begin{enumerate}
  \item Mục tiêu sản phẩm: hệ thống, ứng dụng, mô hình, hay quy trình mà nhóm cần hoàn thành.
  \item Mục tiêu học thuật/kỹ thuật: kiến thức, thuật toán, kiến trúc, hoặc tiêu chí đánh giá cần đạt.
\end{enumerate}

\section{Phạm vi và giả định}

Nên liệt kê rõ những gì đồ án có và không bao gồm. Cách viết này giúp tránh mơ hồ phạm vi ngay từ đầu.

\begin{figure}[H]
  \centering
  \includegraphics[width=0.78\linewidth]{demo-fig.png}
  \caption{Khung giữ chỗ cho sơ đồ tổng quan hệ thống.}
  \label{fig:system-overview}
\end{figure}

Hình \ref{fig:system-overview} là chỗ đặt sẵn để thay bằng sơ đồ kiến trúc, luồng dữ liệu,
hoặc workflow chính của đề tài.
